\section{Conclusion}

This paper documents a clear firm-size wage premium in Colombia. The descriptive evidence shows that real hourly wages rise sharply with firm size, while employment remains concentrated in solo work and micro-firms. Solo workers and firms with up to ten workers account for about two
thirds of employment, while the middle of the firm-size distribution remains thin.

The regression results show that a substantial conditional premium remains after controlling for gender, age, age squared, education, formality, and sector-year fixed effects. In the full specification, workers in firms with 101 or more employees earn about 43 percent more than otherwise comparable solo workers. At the same time, the large reduction in the coefficient as controls are added shows that worker composition, sector-year structure, and formality explain much but not all of the wage gradient.

The formal-informal results sharpen this interpretation. Formal workers earn more than informal workers within every firm-size category, so formality remains a large wage-level margin. However, the conditional wage gradient by firm size is
much steeper among informal workers than among formal workers. Within formal employment, the broad size gradient is weak when formal solo workers are the reference group, although workers in 101+ firms earn more than workers in intermediate-size formal firms. If the firm-size premium is interpreted as a
marker of allocative frictions rather than as pure unobserved sorting, the evidence therefore points to three possible margins for efficiency gains: moving workers and jobs from informality to formality, scaling informal micro-units into larger and more organized firms, and, more selectively, growth or movement within the formal sector from intermediate-size firms toward larger employers.

The central interpretation is therefore deliberately cautious. The firm-size wage premium is a marker of wage inequality and labor-market segmentation, and it is consistent with allocative frictions that limit firm growth or worker movement toward better matches. It is not, by itself, a direct test of misallocation. The GEIH does not identify firms or observe value added, capital, or marginal products, so the paper cannot determine whether the remaining premium reflects unobserved worker sorting or firm-side wage differences, nor whether those differences correspond to allocative inefficiencies. Further research should extend the analysis with firm-level measures of output, capital, employment and productivity to test the allocative-frictions interpretation linked to the household-survey wage patterns documented here.
